\documentclass{article}
\usepackage{a4wide}
\usepackage{amsmath, amsfonts, amssymb}
\usepackage{xcolor}
\usepackage{enumerate}
\input{title.tex}

\newcommand{\N}{\mathbb{N}}
\newcommand{\Z}{\mathbb{Z}}
\newcommand{\R}{\mathbb{R}}
\newcommand{\Q}{\mathbb{Q}}
\newcommand{\F}{\mathbb{F}}
\newcommand{\tr}{\text{tr}}
\newcommand{\Mat}{\text{Mat}}
\newcommand{\ev}{\text{ev}}
\begin{document}
\maketitle
\section{}
\subsection{}
\begin{enumerate}[(a)]
    \item
        \color{blue}
            Sei $m\in\Z$ und $n\geq 2$ eine Natürliche Zahl.
            Ist jede rationale Lösung von $x^n=m$ eine
            Ganze Zahl?\\
        \color{black}
        Es gilt
        $x^n=m\Leftrightarrow x^n-m=0$
        es ist für jede Rationale Lösung
        $x_i=\frac{r_i}{s_i}$ mit $r_i,s_i\in\Z$ Teilerfremd
        nach Theorem 1.2
        $s_i|1\Rightarrow s_i\in\{-1,1\}$
        Also ist $\{x_i\}\subset\Z$.
    \item
        \color{blue}
            Folgern Sie, dass $\sqrt{5}$ und $\sqrt[7]{3}$
            irrational sind.\\
        \color{black}
        \begin{itemize}
            \item
                Wir nehmen an $\sqrt{5}\in\Q$.
                Es gilt $x^2=5$ für $x=\sqrt{5}$.
                Es gilt also nach $(a) \sqrt{5}\in\Z$.
                Da dies Offensichtlich nicht gilt,
                ist die Annahme, $\sqrt{5}\in\Q$, falsch
                $\Rightarrow\sqrt{5}\notin\Q$
            \item
                Wir nehmen an $\sqrt[7]{3}\in\Q$.
                Es gilt $x^7=3$ für $x=\sqrt[7]{3}$.
                Es gilt also nach $(a) \sqrt[7]{3}\in\Z$.
                Da dies Offensichtlich nicht gilt,
                ist die Annahme, $\sqrt[7]{3}\in\Q$, falsch
                $\Rightarrow\sqrt[7]{3}\notin\Q$
        \end{itemize}
    \item
        \color{blue}
        Bestimmtne Sie alle rationalen Nullstellen des
        Polynoms $f(X)=X^{42}+3X^{27}-3X^2-1$\\
        \color{black}
        Nach 1.2 gilt
        $\underset{x_i:=\frac{r_i}{s_i}:f(x_i)=0}\forall
        r_i|-1\Rightarrow r_i\in\{-1,1\},
        s_i|1\Rightarrow s_i\in\{-1,1\}$

        Es ist aber
        $f(-1)=-6,f(1)=0$
        also ist $x_i=\{1\}$
    \item
        \color{blue}
        Bestimmen Sie alle $m\in\Z$, für die $X^3-mX+1$ eine
        rationale Nullstelle besitzt.\\
        \color{black}
        Wir schreiben $f^m(x)=x^3-mx+1$.

        nach 1.2 gilt
        \begin{equation*}
            \underset{x_i^m:=\frac{r_i^m}{s_i^m}:f^m(x_i)=0}
            \forall
            r_i^m|1\Rightarrow r_i^m\in\{-1,1\},
            s_i^m|1\Rightarrow s_i^m\in\{-1,1\}
            \Rightarrow x_i^m\in\{-1,1\}
        \end{equation*}
        Es exisieren also nur die Rationalen
        Nullstellen $-1, 1$
        \begin{itemize}
            \item
                \begin{align*}
                    f^m(-1)&\overset{!}=0\\
                    (-1)^3-m(-1)+1&=0\\
                    -1+m+1&=0\\
                    m&=0
                \end{align*}
            \item
                \begin{align*}
                    f^m(1)&\overset{!}=0\\
                    1^3-m\cdot 1+1&=0\\
                    2&=m
                \end{align*}
        \end{itemize}
        Es exisieren also rationale Nullstellen für
        $f^{0}$ an der Stelle $-1$ und $f^2$ an der Stelle 1.

        Weiter ist $f^0$ das einzige Polynom mit einer
        Nullstelle für $x=-1$ und $f^2$ das einzige Polynom
        mit einer Nullstelle für $x=1$.
        Also ist $m\in\{0,2\}$
\end{enumerate}
\end{document}
